\documentclass{article}

\usepackage{amsmath, amsthm, amssymb, amsfonts}
\usepackage{thmtools}
\usepackage{graphicx}
\usepackage{setspace}
\usepackage{geometry}
\usepackage{float}
\usepackage[hidelinks]{hyperref}
\usepackage[utf8]{inputenc}
\usepackage[spanish]{babel}
\usepackage{framed}
\usepackage[dvipsnames]{xcolor}
\usepackage{tcolorbox}
\usepackage{tikz}
\usepackage{caption}
\usepackage{longtable}
\usepackage{pdflscape}
\usepackage{svg}
\usepackage{subcaption}
\usepackage{caption}
\usepackage{multirow}
\usepackage{array}
\usepackage{listings}
\usepackage{cancel}
\usepackage{fancyhdr}


\colorlet{LightGray}{White!90!Periwinkle}
\colorlet{LightOrange}{Orange!15}
\colorlet{LightGreen}{Green!15}



\newcommand{\HRule}[1]{\rule{\linewidth}{#1}}

\declaretheoremstyle[name=Theorem,]{thmsty}
\declaretheorem[style=thmsty,numberwithin=section]{theorem}
\tcolorboxenvironment{theorem}{colback=LightGray}

\declaretheoremstyle[name=Proposition,]{prosty}
\declaretheorem[style=prosty,numberlike=theorem]{proposition}
\tcolorboxenvironment{proposition}{colback=LightOrange}

\declaretheoremstyle[name=Principle,]{prcpsty}
\declaretheorem[style=prcpsty,numberlike=theorem]{principle}
\tcolorboxenvironment{principle}{colback=LightGreen}

\newcolumntype{L}[1]{>{\raggedleft\let\newline\\\arraybackslash\hspace{0pt}}m{#1}}
\newcolumntype{C}[1]{>{\centering\let\newline\\\arraybackslash\hspace{0pt}}m{#1}}
\newcolumntype{R}[1]{>{\raggedright\let\newline\\\arraybackslash\hspace{0pt}}m{#1}}

\setstretch{1.2}
\geometry{
    textheight=9in,
    textwidth=5.5in,
    top=1in,
    headheight=12pt,
    headsep=25pt,
    footskip=30pt
}

\lstdefinestyle{bashstyle}{
    language=bash,
    basicstyle=\ttfamily,
    backgroundcolor=\color{gray!10},
    keywordstyle=\color{blue},
    commentstyle=\color{green!40!black},
    stringstyle=\color{red},
    showstringspaces=false,
    numbers=left,
    numberstyle=\tiny\color{gray},
    breaklines=true,
    breakatwhitespace=true,
    frame=tb,
    rulecolor=\color{black!70},
    framerule=0.5pt,
    tabsize=4,
    captionpos=b
}

\lstdefinestyle{javastyle}{
    language=Java,
    basicstyle=\ttfamily,
    backgroundcolor=\color{gray!10},
    keywordstyle=\color{blue},
    commentstyle=\color{green!40!black},
    stringstyle=\color{red},
    showstringspaces=false,
    numbers=left,
    numberstyle=\tiny\color{gray},
    breaklines=true,
    breakatwhitespace=true,
    frame=tb,
    rulecolor=\color{black!70},
    framerule=0.5pt,
    tabsize=4,
    captionpos=b
}

\lstdefinestyle{pythonstyle}{
    language=Python,
    basicstyle=\ttfamily,
    backgroundcolor=\color{gray!10},
    keywordstyle=\color{blue},
    commentstyle=\color{green!40!black},
    stringstyle=\color{red},
    showstringspaces=false,
    numbers=left,
    numberstyle=\tiny\color{gray},
    breaklines=true,
    breakatwhitespace=true,
    frame=tb,
    rulecolor=\color{black!70},
    framerule=0.5pt,
    tabsize=4,
    captionpos=b
}

\lstdefinestyle{cppstyle}{
    language=C++,
    basicstyle=\ttfamily,
    backgroundcolor=\color{gray!10},
    keywordstyle=\color{blue},
    commentstyle=\color{green!40!black},
    stringstyle=\color{red},
    showstringspaces=false,
    numbers=left,
    numberstyle=\tiny\color{gray},
    breaklines=true,
    breakatwhitespace=true,
    frame=tb,
    rulecolor=\color{black!70},
    framerule=0.5pt,
    tabsize=4,
    captionpos=b
}



\begin{document}
% ------------------------------------------------------------------------------

\section{Template General}

\begin{lstlisting}[style=cppstyle]
#include <bits/stdc++.h>
using namespace std;
#define fast ios::sync_with_stdio(false); cin.tie(nullptr); cout.tie(nullptr)
#define endl "\n"
#define pb push_back
#define all(x) (x).begin(), (x).end()
#define sz(x) (int)(x).size()


typedef long long ll;
typedef short int si;
typedef unsigned long long ull;
typedef long double ld;
typedef unsigned int ui;
typedef string str;
typedef pair<int, int> pii;
typedef pair<ll, ll> pll;

const int INFI = INT_MAX;
const long INFL = LONG_MAX;
const long long INFLL = LLONG_MAX;



int main(){
    fast;

    si n; cin >> n;
    vector<int> v;
    v.resize(n);
    for (int i = 0; i < n; ++i) cin >> v[i];

    cout << fixed << setprecision(10);

    return 0;
}

\end{lstlisting}


\section{Binary Exponentiation}

\begin{lstlisting}[style=cppstyle]
  ll binpow(ll a, ll b){
    ll res = 1;
    ll counter = 0;
    while (b > 0){
      bitset<16> binary(b);
      if ( b& 1 ) res = res * a;
      a + a * a;
      counter++;
      b >>=1;
    }
    return res;
  }

  ll binpowmod(ll a, ll b, ll m){
    a \%= m;
    ll res = 1;
    while ( b > 0 ){
      if ( b & 1 ) res = res * a \% m;
      a = a * a \% m;
      b >>= 1;
    }
    return res;
  }
\end{lstlisting}

\section{Swap dos números}


\begin{lstlisting}[style=cppstyle]
void swap_two_iters(vector<int> vec, int n){
    const auto itr = find(all(vec), n);
    auto indx = distance(vec.begin(), itr);
    iter_swap(vec.begin() + n, itr);
    iter_swap(vec.begin() + n, vec.begin() + indx);
}
\end{lstlisting}


\section{Bits}

\begin{lstlisting}[style=cppstyle]
  bitset<8> b1; // 00000000
  bitset<8> b2(5); // 00000101
  bitset<8> b3(string("1101")); // 00001101
  bitset<8> b("10110010")
  b[0] // 0 (least significant bit)
  b[7] // 1 (most significant bit)
  b[0] = 1; // Set bit 0 10110011
  b.set(3); // set bit at index 3 10111011
  b.reset(1); // reset bit at index 1
  b.flip(2); // toggle bit at index 2 10111101
  b.flip();

  b.any() // True if any bit == 1
  b.none() // true if all bits == 0
  b.all() // true if all bits == 1
  b.count() // number of bits == 1
  b.size() // number of bits in bitset

  b.to_ulon() // unsigned long
  b.to_ullong() // unsigned long long
  b.to_string() // string
\end{lstlisting}










\end{document}
